\documentclass[11pt,a4paper]{article}

\usepackage[T1]{fontenc}
\usepackage[utf8]{inputenc}
\usepackage[a4paper,margin=2.4cm]{geometry}
\usepackage{xcolor}
\usepackage{titlesec}
\usepackage{enumitem}
\usepackage{fancyhdr}
\usepackage{booktabs}
\usepackage{tcolorbox}
\usepackage{hyperref}

\definecolor{laranja}{HTML}{E35A1C}
\definecolor{grafite}{HTML}{23282F}
\definecolor{cinza}{HTML}{6A7280}
\definecolor{fundo}{HTML}{F4F5F7}

\hypersetup{
  colorlinks=true,
  linkcolor=laranja,
  urlcolor=laranja,
  pdftitle={Remoto GF - Documentacao},
  pdfauthor={Gabriel Francisco Fulgoni}
}

\titleformat{\section}{\Large\bfseries\color{grafite}}{\thesection.}{0.6em}{}
\titleformat{\subsection}{\large\bfseries\color{grafite}}{\thesubsection}{0.6em}{}
\titlespacing*{\section}{0pt}{1.4em}{0.6em}

\pagestyle{fancy}
\fancyhf{}
\lhead{\color{cinza}\small Remoto GF}
\rhead{\color{cinza}\small Documentacao do projeto}
\cfoot{\color{cinza}\small \thepage}
\renewcommand{\headrulewidth}{0.4pt}

\newtcolorbox{nota}[1][]{colback=fundo,colframe=laranja,boxrule=0.8pt,
  arc=3pt,left=8pt,right=8pt,top=6pt,bottom=6pt,#1}

\newcommand{\feito}{\textcolor{laranja}{\textbf{[Pronto]}}\ }
\newcommand{\pendente}{\textcolor{cinza}{\textbf{[A fazer]}}\ }

\setlist{itemsep=2pt,topsep=4pt}

\begin{document}

\begin{titlepage}
\centering
\vspace*{3cm}
{\Huge\bfseries\color{laranja} Remoto GF\par}
\vspace{0.6cm}
{\LARGE\color{grafite} Documentacao do projeto\par}
\vspace{0.4cm}
{\large\color{cinza} Controle do PC pelo celular, e o plano de controlar o celular pelo PC\par}
\vspace{2cm}
\begin{tcolorbox}[colback=fundo,colframe=laranja,width=0.8\textwidth,halign=center,arc=4pt]
\large Autor: \textbf{Gabriel Francisco Fulgoni}\\[4pt]
Editor de video e pos-producao audiovisual\\[4pt]
Repositorio: \url{https://github.com/GhunterGh931/remoto-gf}
\end{tcolorbox}
\vfill
{\color{cinza} Barueri, SP \quad\textbullet\quad Versao 1.0 \quad\textbullet\quad 19 de outubro de 2026\par}
\end{titlepage}

\tableofcontents
\newpage

% ============================================================
\section{Visao geral}

O \textbf{Remoto GF} e um projeto pessoal que transforma o celular em um controle
remoto completo do computador, pela rede Wi-Fi local. O celular abre uma pagina
no navegador (ou um aplicativo instalado) e comanda o PC: mouse, teclado, volume,
aplicativos abertos, tela ao vivo e mais. Nao e preciso cadastro, nuvem ou internet
externa: tudo acontece entre o celular e o PC na mesma rede.

O projeto tem duas direcoes de controle:

\begin{itemize}
  \item \feito \textbf{Celular controla o PC.} Concluido e funcionando.
  \item \feito \textbf{PC controla o celular e troca arquivos.} Concluido e funcionando
  dentro do proprio Remoto GF (ver Secao~\ref{sec:futuro}).
\end{itemize}

\begin{nota}
\textbf{Principio de seguranca do projeto:} tudo roda apenas na rede local do usuario,
o acesso e protegido por um PIN, e nenhuma acao sensivel acontece sem a confirmacao
da pessoa. O projeto nao coleta nem envia dados para fora.
\end{nota}

% ============================================================
\section{Arquitetura}

\subsection{Como as pecas se conectam}

\begin{enumerate}
  \item \textbf{Servidor no PC} (\texttt{remoto.py}, em Python): abre um servidor web
  e WebSocket na porta 8765, exibe uma janela com QR Code e PIN, e executa no PC os
  comandos que chegam do celular.
  \item \textbf{Interface do celular} (\texttt{web/index.html}): uma unica pagina com
  todos os paineis de controle. Funciona tanto no navegador do celular quanto dentro
  do aplicativo Android.
  \item \textbf{Aplicativo Android} (pasta \texttt{android/}): empacota a mesma interface
  em um app instalavel e adiciona recursos que so um app tem (botoes de volume,
  notificacoes do sistema, giroscopio, busca automatica do PC na rede).
\end{enumerate}

\subsection{Conexao e pareamento}

\begin{itemize}
  \item Em casa, celular e PC ficam na mesma rede (Wi-Fi ou cabo no PC).
  \item Fora de casa (dados moveis), a conexao passa pelo Tailscale (ver
  Secao~\ref{sec:fora}).
  \item O PC mostra um QR Code e um PIN de 6 digitos.
  \item O celular escaneia o QR (ou digita o endereco e o PIN) e fica pareado por um
  token salvo no aparelho.
  \item No aplicativo, o celular encontra o PC sozinho na rede (busca por UDP broadcast).
  \item Protecao: apos 5 erros de PIN, o pareamento fica bloqueado por 1 minuto, e o PC
  avisa a tentativa.
\end{itemize}

% ============================================================
\section{O que ja foi feito: celular controla o PC}
\label{sec:feito}

Todos os itens abaixo estao \feito implementados e testados.

\subsection{Touchpad e mouse}
\begin{itemize}
  \item Area de toque com cursor, 1 dedo move, toque clica.
  \item Toque com 2 dedos: clique direito. Toque com 3 dedos: clique do meio.
  \item Pressao longa: clique direito. Toque duplo e arrastar: selecionar e mover.
  \item Rolagem com 2 dedos e faixa lateral de rolagem.
  \item Pinca para zoom (Ctrl + roda do mouse). Arrastar com 3 dedos.
  \item Botoes fisicos do mouse (esquerdo, meio, direito) e trava de clique para arrastar.
  \item Botao giroscopio: mover o cursor inclinando o celular.
\end{itemize}

\subsection{Teclado}
\begin{itemize}
  \item Digitacao ao vivo: o que se digita aparece no PC na hora.
  \item Teclas modificadoras (Ctrl, Alt, Shift, Win), setas, Enter, Esc, Tab, Backspace.
  \item Teclas de funcao (F1 a F12) e teclado numerico.
  \item Modo buffer (digitar e enviar com Enter) e modificador por toque.
\end{itemize}

\subsection{Atalhos}
\begin{itemize}
  \item Atalhos gerais do Windows (copiar, colar, desfazer, Alt+Tab, print, etc).
  \item Atalhos do \textbf{DaVinci Resolve} (play/pause, J K L, marcadores, corte, lamina,
  zoom, etc), feitos sob medida para o fluxo de edicao.
  \item Atalhos personalizados que o usuario cria e salva.
\end{itemize}

\subsection{Controle de midia e som}
\begin{itemize}
  \item Volume geral do PC, com mudo e passos de 5\%.
  \item Volume por aplicativo (mixer): controlar cada programa que esta tocando som.
  \item Controles de midia: play/pause, anterior, proximo, parar.
  \item Brilho da tela (quando o PC permite).
\end{itemize}

\subsection{Aplicativos e janelas}
\begin{itemize}
  \item Lista das janelas abertas, com focar, minimizar, maximizar, fechar e forcar fim.
  \item Abrir aplicativos a partir de uma lista configuravel.
  \item Abrir sites ou pesquisas no navegador do PC.
\end{itemize}

\subsection{Tela ao vivo}
\begin{itemize}
  \item Ver a tela do PC no celular, com escolha de monitor, taxa de quadros e qualidade.
  \item Tocar na imagem clica naquele ponto; segurar faz clique direito.
  \item Correcao: antes, todo toque virava clique direito no app Android. O sistema
  cancelava o toque na imagem e o fim do toque nao chegava, entao o cronometro do clique
  direito disparava. Agora o app usa eventos de toque e qualquer cancelamento desarma o
  clique direito.
  \item Tela deitada com mouse: o botao no canto da imagem abre a tela na horizontal com
  70\% para a tela do PC e 30\% para um touchpad (mesmos gestos do painel Mouse) com botoes
  Esquerdo e Direito. ``So tela'' esconde o touchpad e ``Sair'' (ou o botao Voltar) fecha.
\end{itemize}

\subsection{Energia e sistema}
\begin{itemize}
  \item Bloquear, desligar a tela, suspender, reiniciar e desligar (com cancelamento).
  \item Desligamento agendado por tempo.
  \item Status do PC: CPU, memoria, bateria e tempo ligado.
  \item Copiar e colar texto entre o celular e o PC.
\end{itemize}

\subsection{Configuracoes e experiencia}
\begin{itemize}
  \item Ordem dos paineis arrastavel e paineis que podem ser escondidos.
  \item Controle com uma mao, modo canhoto, velocidades de movimento e rolagem.
  \item Aceleracao, precisao, reducao de latencia, feedback por som e vibracao.
  \item Reconexao automatica, tela sempre ligada, bloqueio de orientacao.
\end{itemize}

\subsection{Avisos de autorizacao no celular}
O app avisa, por notificacao, quando: um programa esta aberto como administrador
(o Windows bloqueia o controle remoto nele), aparece a janela de permissao do Windows,
alguem erra o PIN varias vezes, um novo aparelho conecta, ou um desligamento foi agendado.

\subsection{Distribuicao}
\begin{itemize}
  \item Repositorio publico no GitHub.
  \item Compilacao automatica (GitHub Actions): a cada atualizacao do codigo, o servico
  gera o \texttt{RemotoGF.apk} (celular) e o \texttt{RemotoGF-Instalador.exe} (PC) e publica
  na pagina de Releases.
  \item Scripts locais: \texttt{iniciar.bat} (roda sem instalar) e \texttt{gerar\_exe.bat}.
\end{itemize}

\begin{nota}
\textbf{Links de download:}\\
Celular (APK): \url{https://github.com/GhunterGh931/remoto-gf/releases/latest}\\
PC (instalador): mesma pagina de Releases.
\end{nota}

% ============================================================
\section{PC controla o celular e transferencia de arquivos}
\label{sec:futuro}

\subsection{Recursos implementados}
Controlar o celular a partir do computador, pela rede Wi-Fi local, diretamente dentro do
proprio aplicativo Remoto GF (no modo ``Deixar o PC me controlar'') e trocar arquivos nos
dois sentidos:

\begin{itemize}
  \item \feito \textbf{Ver a tela do celular no PC e controlar com mouse e teclado:}
  clique toca, arraste ou rodinha do mouse desliza, segurar faz pressao longa, teclado
  digita ao vivo (incluindo \texttt{Ctrl+V} para colar texto do PC direto no celular).
  \item \feito \textbf{Abrir e usar qualquer aplicativo do celular pelo PC, com foco no WhatsApp:}
  botao dedicado ``WhatsApp'' na janela do PC e menu ``Abrir app'' que lista e abre
  qualquer aplicativo instalado no celular remotamente.
  \item \feito \textbf{Passar arquivos entre o PC e o celular, nos dois sentidos:}
  aba ``Arquivos'' no celular para enviar arquivos ao PC (salvos em \texttt{Downloads\textbackslash Remoto GF})
  e botao ``Enviar arquivo ao celular'' na janela do PC para disponibilizar arquivos que
  o celular baixa com um toque.
\end{itemize}

\subsection{Como funciona no aplicativo}
\begin{enumerate}
  \item \textbf{Acessibilidade (uma unica vez):} nas configuracoes do Android, o usuario ativa
  o servico \texttt{Remoto GF: controle pelo PC} (\texttt{ControlService}), que permite executar
  toques, gestos e digitacao enviados pelo PC pareado.
  \item \textbf{Transmissao de tela (\texttt{ScreenService}):} ao tocar em \textbf{Iniciar} na aba
  ``PC controla'', o Android pede confirmacao de captura de tela (\texttt{MediaProjection}).
  Enquanto a sessao esta ativa, uma notificacao fixa permanece visivel com o botao \textbf{Parar}.
  \item \textbf{Janela no PC:} a janela do celular abre automaticamente no PC (ou pelo botao
  ``Ver meu celular''), com barra de navegacao (Voltar, Inicio, Recentes, Bloquear, Print),
  seletor de qualidade, atalho direto para o WhatsApp, menu de aplicativos do celular e envio
  de arquivos.
\end{enumerate}

\begin{nota}
\textbf{Observacao tecnica:} o Bluetooth nao tem banda para transmitir video de tela com
boa qualidade. Por isso o espelhamento e a transferencia de arquivos acontecem pelo Wi-Fi
local, protegidos pelo mesmo pareamento por PIN.
\end{nota}

% ============================================================
\section{Downloads do PC no celular}

O Baixador Drive GF grava o progresso a cada segundo em
\texttt{\%LOCALAPPDATA\%\textbackslash GF\textbackslash progresso\textbackslash baixador-drive-gf.json}.
O Remoto GF le essa pasta, ignora arquivos velhos (app fechado ou travado) e manda as mudancas
ao celular (mensagem \texttt{progress}). No celular:
\begin{itemize}
  \item Notificacao fixa com barra de progresso, sem som (canal ``Downloads no PC''): arquivo,
  porcentagem, velocidade, tempo que falta e posicao na fila. Ao concluir, parar ou dar erro,
  ela vira um aviso com som.
  \item Faixa no topo do app; tocando nela aparecem os detalhes do arquivo e da fila.
\end{itemize}
Qualquer app pode aparecer gravando um JSON na mesma pasta, com os campos \texttt{estado}
(baixando, pausado, concluido, parado, erro), \texttt{arquivo}, \texttt{pct}, \texttt{feito},
\texttt{total}, \texttt{vel}, \texttt{resta}, \texttt{geral} e \texttt{msg}. Limite: a
notificacao so atualiza enquanto o app do celular esta aberto ou em segundo plano.

% ============================================================
\section{Acesso fora de casa (dados moveis)}
\label{sec:fora}

\subsection{O problema}
Pelos dados moveis o celular nao enxerga o IP local do PC (192.168.x.x). Abrir porta no
roteador costuma falhar, porque muitas operadoras usam CGNAT, e ainda deixaria o PC exposto
na internet.

\subsection{A solucao: Tailscale}
O Tailscale cria uma rede privada e criptografada entre os aparelhos da mesma conta. O PC
ganha um IP fixo na faixa 100.x.x.x, que funciona de qualquer lugar. Nada fica aberto para a
internet: so os aparelhos logados na conta enxergam o PC, e o Remoto GF ainda exige o token
do pareamento. Recomendado: senha forte e verificacao em duas etapas na conta do Tailscale.

\subsection{Como usar}
\begin{enumerate}
  \item Instale o Tailscale no PC e no celular e entre com a mesma conta nos dois.
  \item Na janela do Remoto GF no PC aparece a linha ``Fora de casa (Tailscale)'' com o
  endereco 100.x.x.x:8765.
  \item Conecte o celular em casa uma vez: o app aprende o endereco fora de casa sozinho.
  Ele tambem pode ser digitado na tela de pareamento ou em Configuracoes, secao
  ``Fora de casa''.
  \item Fora de casa, deixe o Tailscale ligado no celular e abra o Remoto GF.
\end{enumerate}

\subsection{Como o app escolhe o endereco}
O app guarda dois enderecos: o de casa e o de fora de casa. Ele tenta o endereco usado por
ultimo; se nao abrir em alguns segundos, troca para o outro sozinho e guarda o que
funcionou. O Firewall do Windows ja libera o Remoto GF em qualquer rede (perfil ``any'' no
instalador), entao a rede do Tailscale passa sem configuracao extra.

\subsection{Wake-on-LAN (ligar o PC pela rede)}
O pacote magico so funciona dentro da rede local, entao fora de casa ele precisa de um
aparelho sempre ligado na rede de casa para envia-lo (planejado para depois). Correcao feita:
o PC agora informa o MAC da placa que tem o IP local (cabo ou Wi-Fi). Antes, ele podia
informar a placa de um adaptador virtual (VirtualBox, VPN), e o pacote ia para o endereco
errado. Requisitos no PC: Wake-on-LAN ativo na BIOS, ``Wake on Magic Packet'' ativo na placa
de rede e inicializacao rapida do Windows desligada. Depois de atualizar o PC, conecte o
celular uma vez em casa para ele salvar o MAC correto.

% ============================================================
\section{Resumo do estado}

\begin{center}
\begin{tabular}{@{}p{8.5cm}l@{}}
\toprule
\textbf{Item} & \textbf{Estado} \\
\midrule
Celular controla o PC (todas as funcoes) & Pronto \\
Aplicativo Android instalavel & Pronto \\
Instalador do PC & Pronto \\
Compilacao automatica no GitHub & Pronto \\
Repositorio publico & Pronto \\
PC ve e controla o celular (com atalho WhatsApp e apps) & Pronto \\
Passar arquivos entre PC e celular (nos dois sentidos) & Pronto \\
Acesso fora de casa pelos dados moveis (Tailscale) & Pronto \\
Ligar o PC pela rede (Wake-on-LAN) em casa & MAC corrigido, testar \\
Ligar o PC pela rede fora de casa & Planejado \\
Progresso dos downloads do PC no celular & Pronto \\
\bottomrule
\end{tabular}
\end{center}

% ============================================================
\section{Proximos passos sugeridos}

\begin{enumerate}
  \item Escolher entre Vincular ao Celular e scrcpy para o controle do celular pelo PC.
  \item Configurar a ferramenta escolhida (passo a passo pode ser feito em seguida).
  \item Opcional: no Remoto GF, proteger a chave de assinatura do app movendo-a para os
  Secrets do GitHub, caso o projeto seja distribuido de forma mais ampla.
  \item Opcional: adicionar transferencia de arquivos pela rede no Remoto GF (do PC para o
  celular e do celular para o PC), que nao envolve controle remoto e e viavel.
\end{enumerate}

\vfill
\begin{center}
\color{cinza}\small Documento gerado para registro do projeto Remoto GF.
\end{center}

\end{document}
